\documentclass[10pt,a4paper]{amsart}
\usepackage[margin=19mm]{geometry}
\usepackage{amsmath,amssymb,mathtools}
\usepackage[scr=rsfs,cal=euler]{mathalfa}
\usepackage{xfrac}
\usepackage[unicode,hidelinks,bookmarks=false]{hyperref}
\newcommand{\ie}{i.e.\ }
\newcommand{\cf}{cf.\ }
\newcommand{\resp}{respectively\ }
\newcommand{\Subsection}{\subsection}
\providecommand{\nobreakdash}{\nobreak\hbox{-}}
\setlength{\emergencystretch}{2em}
\multlinegap0pt
\usepackage{siunitx}
\usepackage{siunitx}
\usepackage{siunitx}
\usepackage{siunitx}
\usepackage{amssymb}
\usepackage{algorithm}\usepackage{algpseudocode}
\usepackage{siunitx}
\usepackage{xcolor}
\usepackage{dsfont}
\usepackage{bm}
\newtheorem{assumption}{Assumption}
\newtheorem{proposition}{Proposition}
\title{Precision autotuning for linear solvers via contextual bandit-based RL}
\author{Erin Carson, Xinye Chen}
\date{Source: arXiv:2601.00728v5 (CC BY 4.0)}
\begin{document}
\maketitle

\noindent\emph{Source and licence.} Precision autotuning for linear solvers via contextual bandit-based RL. Version arXiv:2601.00728v5 (CC BY 4.0), distributed by arXiv under the Creative Commons Attribution 4.0 International licence, http://creativecommons.org/licenses/by/4.0/, as recorded in the arXiv licence metadata for this exact version and shown on the abstract page https://arxiv.org/abs/2601.00728v5. Full licence notice: \texttt{licenses/arxiv-2601.00728-CC-BY-4.0.txt}.\par\smallskip

\noindent\textit{Editorial scope.} Counted unit: the article's proposition prop:discretization (source0.tex lines 325--347). The article's complete original proof of it is delivered with it (source0.tex lines 349--377, one \texttt{proof} environment of 29 lines). No mathematics is written, repaired or re-derived: the statement and the proof are the article's own lines, in the article's own notation, and no retained line is substituted or re-flowed.  Retained uncounted as the source-local context the proof needs: lines 313--321.  Nothing was omitted from any retained span, so the package carries no omission record.  No retained line contains a citation, so the wrapper carries no bibliography and no bibliography entry.  No retained line contains a figure environment, an image-inclusion command or drawing code, so no image is delivered and the package declares no asset dependency.  Editorial formatting only, in addition: an AMS \texttt{amsart} preamble that restates the article's own declarations and macros used by the retained lines, namely the environments \texttt{assumption}, \texttt{proposition} on separate counters, together with the standard packages those lines need.  The retained environments print in this document as Assumption 1, Proposition 1 in order of appearance, whereas the article prints the same items with the numbering of its own full document.  The complete native source of the article as distributed in the arXiv e-print for this exact version is bundled unchanged as \texttt{sources/192020/source0.tex}.
\begin{assumption}[Lipschitz continuity of the expected reward]
\label{ass:lipschitz}
Let $\mu(s,\mathbf{a}) := \mathbb{E}[r(s,\mathbf{a})]$ denote the expected reward.
We assume, for all $\mathbf{a} \in \mathcal{A}$, there exists a constant $L > 0$ such that
\begin{equation}
\label{eq:lipschitz}
|\mu(s,\mathbf{a}) - \mu(s',\mathbf{a})| \le L \| s - s' \|, \quad \forall s, s' \in \mathcal{S}.
\end{equation}
\end{assumption}

\begin{proposition}[Discretization error bound]
\label{prop:discretization}
Let $s \in \mathcal{S}$ be a continuous context, and let $s_d \in \mathcal{S}_d$ be its discretized representation.
Assume that each discretization bin has diameter at most $\Delta$, i.e.,
\[
\| s - \omega(s_d) \| \le \Delta,
\]
where $\omega(s_d)$ denotes a representative point (e.g., bin center) of $s_d$.

Define the optimal action in the continuous space and the discretized policy, respectively, as
\begin{equation*}
\begin{aligned}
\mathbf{a}^*(s) &= \arg\max_{\mathbf{a} \in \mathcal{A}} \mu(s,\mathbf{a}), \\ \mathbf{a}_d^*(s_d) &= \arg\max_{\mathbf{a} \in \mathcal{A}} \mu(\omega(s_d), \mathbf{a}).
\end{aligned}
\end{equation*}


Under Assumption~\ref{ass:lipschitz}, the performance loss due to discretization is bounded by
\begin{equation}
\label{eq:discretization_bound}
\mu(s, \mathbf{a}^*(s)) - \mu(s, \mathbf{a}_d^*(s_d)) \le 2 L \Delta.
\end{equation}
\end{proposition}

\begin{proof}
Let $\omega = \omega(s_d)$ be the representative point associated with the discretized state $s_d$.
By Assumption~\ref{ass:lipschitz}, for any $\mathbf{a} \in \mathcal{A}$, we have
\begin{equation}
\label{eq:lipschitz_step}
|\mu(s,\mathbf{a}) - \mu(\omega,\mathbf{a})| \le L \| s - \omega \| \le L \Delta.
\end{equation}

By definition of $\mathbf{a}^*(s)$,
\[
\mu(s,\mathbf{a}^*(s)) \le \mu(\omega,\mathbf{a}^*(s)) + L\Delta.
\]

Since $\mathbf{a}_d^*(s_d)$ maximizes $\mu(\omega,\mathbf{a})$ over $\mathcal{A}$, we have
\[
\mu(\omega,\mathbf{a}^*(s)) \le \mu(\omega,\mathbf{a}_d^*(s_d)).
\]

Applying the Lipschitz bound again,
\[
\mu(\omega,\mathbf{a}_d^*(s_d)) \le \mu(s,\mathbf{a}_d^*(s_d)) + L\Delta.
\]

Combining the above inequalities yields
\[
\mu(s,\mathbf{a}^*(s)) \le \mu(s,\mathbf{a}_d^*(s_d)) + 2L\Delta,
\]
which proves the result.
\end{proof}

\end{document}
